PythonTutor\cite{pythontutor} is a widely used educational tool for visualizing program execution, originally developed to help novices build accurate mental models of how code runs. Its core contribution is a step-by-step execution trace---illustrated in Fig.~\ref{fig:pythontutor}---that exposes stack frames, heap objects, and control flow in a concrete, visual manner. As described by its creator, Philip J. Guo, PythonTutor aims to ``help people overcome a fundamental barrier to learning programming: understanding what happens as the computer runs each line of code'' \cite{Guo2013}.

\begin{figure*}[htbp]
\centerline{\includegraphics[width=\textwidth]{PythonTutor.PNG}}
\caption{Example \href{https://pythontutor.com}{PythonTutor.com} output showing a step-by-step execution trace with stack frames and heap objects.}
\label{fig:pythontutor}
\end{figure*}

PythonTutor has been shown to be effective in improving students' understanding of programming concepts; Guo's further studies demonstrated significant learning gains when PythonTutor was used as a supplement to traditional instruction \cite{Guo2015}. This demonstrates the educational value of connecting code execution to visual representations, as well as the benefits of an accessible, interactive tool.

However, the visualizations produced by PythonTutor are relatively basic and not suitable for presentation purposes. PythonTutor's primary focus is on clarity and educational value rather than aesthetics, resulting in output that may be difficult to read or visually unappealing when projected to an audience. This, along with the fact that the tool is web-based and not designed for offline use, makes it less than ideal for use in settings such as conference presentations where visual impact is important. PythonTutor is additionally restricted by its limited customization options, minimal library of visualization methods, and support for only C, C\thinspace\verb|++|, Java, Javascript, and Python.
